\section{One-anchor rigidity}\label{sec:anchor}

We isolate the combinatorial statement used on the top block. Let $Y\to\infty$, and let $Q$ be a set of primes in $[Y,2Y)$ satisfying
\[
c\frac{Y}{\log Y}
\le N:=|Q|
\le C\frac{Y}{\log Y}.
\tag{4.1}\label{eq:4.1}
\]
An assignment is a tuple $a=(a_p)_{p\in Q}$ with $a_p\in\Z/p\Z$. For distinct $p,q\in Q$, let $H_{pq}(a)$ be the centred simultaneous CRT lift
\[
H_{pq}\equiv a_p\pmod p,
\qquad
H_{pq}\equiv a_q\pmod q,
\qquad
-\frac{pq}{2}<H_{pq}\le\frac{pq}{2}.
\tag{4.2}\label{eq:4.2}
\]
Define
\[
\mathcal Q(a)
=
\sum_{p<q\in Q}
\left(\frac{H_{pq}(a)}{pq}\right)^2,
\qquad
\sigma_Q^2
=
\sum_{p<q\in Q}\frac1{p^2q^2}.
\tag{4.3}\label{eq:4.3}
\]
Then
\[
\sigma_Q^2\asymp\frac1{Y^2\log^2Y}.
\tag{4.4}\label{eq:4.4}
\]
The main conclusion of this section is that energy below the scale $Y/\log^3Y$ forces an exact integer state on every coordinate. The logic has two complementary parts. Low energy first forces one dominant integer label, and reciprocal dispersion then eliminates every exceptional coordinate. Above the rigidity threshold, a fingerprint determines almost all coordinates; the entropy of its completions is small enough to be absorbed by the Bernoulli damping. Thus the first mechanism identifies the coherent states, while the second makes the remaining top assignments negligible in weighted sum.

\subsection{Reciprocal dispersion}

\begin{lemma}[Reciprocal dispersion]\label{lem:reciprocal-dispersion}
Let $q$ be prime with $Y\le q<4Y$, let $F$ be a set of $s\ge16$ primes in $[Y,2Y)$, none equal to $q$, and let $d\not\equiv0\pmod q$. Then at least $s/2$ primes $p\in F$ satisfy
\[
\left\|\frac{d p^{-1}}q\right\|
\ge \frac{s}{64Y}.
\]
Consequently
\[
\sum_{p\in F}
\left\|\frac{d p^{-1}}q\right\|^2
\ge c_D\frac{s^3}{Y^2}
\tag{4.5}\label{eq:4.5}
\]
for an absolute constant $c_D>0$.
\end{lemma}

\begin{proof}
Put $\delta=s/(64Y)$. If
\[
\left\|\frac{dp^{-1}}q\right\|\le\delta,
\]
choose the nearest integer $\ell$ and write
\[
v=dp^{-1}-\ell q.
\]
Then
\[
|v|\le\delta q\le s/16,
\qquad
vp\equiv d\pmod q.
\]
For a fixed nonzero integer $v$, the last congruence places $p$ in one residue class modulo $q$. Since $[Y,2Y)$ has length at most $q$, such a class contains at most two integers. There are at most $2(s/16)+1$ possible values of $v$. Hence fewer than $s/2$ primes can satisfy the displayed small-distance condition. At least $s/2$ primes therefore have distance at least $\delta$, and summing their squares proves \eqref{eq:4.5}.
\end{proof}

The hypothesis $s\ge16$ will be retained explicitly whenever Lemma \ref{lem:reciprocal-dispersion} is used.

\subsection{Repulsion between label classes}

Fix $\rho\in(0,1/4)$. An integer $m$ is called a \emph{dominant label} for $a$ if $|m|\le Y^2/2$ and
\[
a_p\equiv m\pmod p
\]
for at least $(1-\rho)N$ primes $p\in Q$. Two dominant labels are equal: their classes overlap in at least two primes, whose product divides their difference, while the difference has absolute value at most $Y^2$.

\begin{lemma}[Cross-label energy]\label{lem:cross-label}
Let $C_m,C_{m'}\subseteq Q$ be disjoint label classes with $m\ne m'$, and suppose
\[
|m|,|m'|\le B_0<Y^2/4,
\]
\[
|C_m|\ge\max\left(16,\frac{128B_0}{Y}\right),
\qquad
|C_{m'}|\ge2.
\tag{4.6}\label{eq:4.6}
\]
Then
\[
\sum_{p\in C_m}
\sum_{q\in C_{m'}}
\left(\frac{H_{pq}}{pq}\right)^2
\ge
c_E\frac{|C_m|^3(|C_{m'}|-1)}{Y^2}.
\tag{4.7}\label{eq:4.7}
\]
\end{lemma}

\begin{proof}
Put $d=m'-m$. Since $0<|d|<Y^2/2$, at most one prime in $C_{m'}$ can divide $d$. Fix any other $q\in C_{m'}$. For $p\in C_m$, write
\[
H_{pq}=m+jp.
\]
Reducing modulo $q$ gives $jp\equiv d\pmod q$, and therefore
\[
\left\|\frac{dp^{-1}}q\right\|
\le
\frac{|H_{pq}|}{pq}+\frac{|m|}{pq}.
\tag{4.8}\label{eq:4.8}
\]
By Lemma \ref{lem:reciprocal-dispersion}, for at least half the primes $p\in C_m$ the left side is at least $|C_m|/(64Y)$. The size condition in \eqref{eq:4.6} gives
\[
\frac{|m|}{pq}
\le
\frac{B_0}{Y^2}
\le
\frac{|C_m|}{128Y}.
\]
Thus at least half the $p$ satisfy
\[
\frac{|H_{pq}|}{pq}
\ge
\frac{|C_m|}{128Y}.
\]
The energy contributed by this fixed $q$ is $\gg |C_m|^3/Y^2$. Summing over all but at most one $q\in C_{m'}$ gives \eqref{eq:4.7}.
\end{proof}

The factor $|C_{m'}|-1$ in \eqref{eq:4.7} is important and will be kept when the class energies are summed.

\subsection{Low energy forces a dominant label}

\begin{proposition}[Dominant-label forcing]\label{prop:dominant-label}
There is $c_w>0$ such that, for all sufficiently large $Y$,
\[
\mathcal Q(a)<c_w\frac{Y}{\log^3Y}
\tag{4.9}\label{eq:4.9}
\]
implies that $a$ has a unique dominant label.
\end{proposition}

\begin{proof}
If $\mathcal Q(a)=0$, every centred pair lift vanishes, so every coordinate is the residue of the integer $0$. We may therefore assume
\[
0<R:=\mathcal Q(a)<c\frac{Y}{\log^3Y},
\tag{4.10}\label{eq:4.10}
\]
with $c$ to be chosen sufficiently small.

Set
\[
B_0=A_\rho\frac{\sqrt R\,Y^2}{N},
\tag{4.11}\label{eq:4.11}
\]
where $A_\rho$ is a large fixed constant. Call an edge $\{p,q\}$ bad if $|H_{pq}|>B_0$. Since $pq<4Y^2$, every bad edge contributes more than $B_0^2/(16Y^4)$ to $R$, so
\[
\#\{\text{bad edges}\}
\le
\frac{16RY^4}{B_0^2}.
\tag{4.12}\label{eq:4.12}
\]
Some base prime $p_0$ has bad degree at most
\[
\frac{32RY^4}{B_0^2N}
=
\frac{32N}{A_\rho^2}.
\tag{4.13}\label{eq:4.13}
\]
Choose $A_\rho$ so that this is at most $\rho N/8$.

For every nonbad neighbour $q$ of $p_0$, set
\[
\ell(q)=H_{p_0q}.
\]
Then $|\ell(q)|\le B_0$, $\ell(q)\equiv a_q\pmod q$, and all $\ell(q)$ lie in one residue class modulo $p_0$. For every integer label $m$ that occurs, define its label class by
\[
C_m=\{q:\{p_0,q\}\text{ is nonbad and }\ell(q)=m\}.
\]
These classes partition the nonbad neighbours of $p_0$, and every $q\in C_m$ satisfies $a_q\equiv m\pmod q$. The number $M_0$ of possible integer labels therefore satisfies
\[
M_0\le\frac{2B_0}{Y}+2.
\tag{4.14}\label{eq:4.14}
\]
For sufficiently small $c$ and large $Y$, \eqref{eq:4.10}--\eqref{eq:4.11} also give $B_0<Y^2/4$.

Assume for contradiction that no class is dominant. Call a label class small if its size is less than
\[
s_0=256\left(\frac{B_0}{Y}+1\right).
\tag{4.15}\label{eq:4.15}
\]
Then $s_0\ge16$ and $s_0\ge128B_0/Y$, so Lemma \ref{lem:cross-label} applies to every substantial class.

Suppose first that the small classes contain at least $\rho N/4$ vertices. Then
\[
\frac{\rho N}{4}
\le M_0s_0
\le C_\rho\left(\frac{B_0}{Y}+1\right)^2.
\tag{4.16}\label{eq:4.16}
\]
If $B_0/Y<1$, the right side is bounded while $N\to\infty$, a contradiction. Otherwise \eqref{eq:4.11} and \eqref{eq:4.16} imply
\[
N\le C_\rho\frac{B_0^2}{Y^2}
=C_\rho A_\rho^2\frac{RY^2}{N^2},
\]
and hence
\[
R\gg\frac{N^3}{Y^2}\gg\frac{Y}{\log^3Y},
\]
contradicting \eqref{eq:4.10} for sufficiently small $c$.

We may therefore assume that the small classes contain fewer than $\rho N/4$ vertices. After removing $p_0$ and its bad neighbours, the substantial classes have total size
\[
S\ge(1-\rho/2)N
\tag{4.17}\label{eq:4.17}
\]
for large $Y$. Let their sizes be $n_1,\ldots,n_t$. Each $n_i\ge s_0$, while no $n_i$ exceeds $(1-\rho)N$.

Apply Lemma \ref{lem:cross-label} to ordered pairs of distinct substantial classes. Every undirected cross edge is counted twice, so, after absorbing the factor $1/2$ into the constant,
\[
R
\ge
\frac{c}{Y^2}
\sum_i n_i^3
\sum_{j\ne i}(n_j-1).
\tag{4.18}\label{eq:4.18}
\]
The inner sum is exactly
\[
\sum_{j\ne i}(n_j-1)
=(S-n_i)-(t-1).
\tag{4.19}\label{eq:4.19}
\]
Because every other substantial class has size at least $s_0$,
\[
S-n_i\ge(t-1)s_0,
\]
and hence
\[
(S-n_i)-(t-1)
\ge
\left(1-\frac1{s_0}\right)(S-n_i)
\ge
\frac{255}{256}(S-n_i).
\tag{4.20}\label{eq:4.20}
\]
Thus
\[
R\ge\frac{c}{Y^2}\sum_i n_i^3(S-n_i).
\tag{4.21}\label{eq:4.21}
\]

If the largest class has size at least $S/2$, then it contains $\gg N$ vertices while \eqref{eq:4.17} and non-dominance leave $\gg_\rho N$ substantial vertices outside it. Its single term in \eqref{eq:4.21} is therefore $\gg_\rho N^4/Y^2$, too large for \eqref{eq:4.10}.

If every class has size less than $S/2$, then $S-n_i\ge S/2$, and power mean gives
\[
\sum_i n_i^3
\ge
\frac{S^3}{t^2}
\ge
\frac{S^3}{M_0^2}.
\]
Hence
\[
R\ge c_\rho\frac{N^4}{M_0^2Y^2}.
\tag{4.22}\label{eq:4.22}
\]
If $B_0/Y<1$, this again exceeds the scale in \eqref{eq:4.10}. Otherwise, by \eqref{eq:4.14},
\[
M_0^2
\le C\frac{B_0^2}{Y^2}
=C A_\rho^2\frac{RY^2}{N^2}.
\]
Substitution into \eqref{eq:4.22} yields
\[
R^2\ge c\frac{N^6}{Y^4},
\qquad
R\ge c\frac{N^3}{Y^2}\gg\frac{Y}{\log^3Y},
\]
again a contradiction. Choosing $c_w$ below the finitely many constants above proves the proposition.
\end{proof}

\subsection{No exceptional coordinates}

\begin{proposition}[Exact zero-exception rigidity]\label{prop:zero-exception}
After decreasing $c_w$ if necessary, every assignment satisfying \eqref{eq:4.9} is represented by one integer $m$:
\[
a_p\equiv m\pmod p
\qquad(p\in Q).
\tag{4.23}\label{eq:4.23}
\]
Moreover
\[
|m|
\le
C_\rho\frac{\sqrt{\mathcal Q(a)}}{\sigma_Q}.
\tag{4.24}\label{eq:4.24}
\]
\end{proposition}

\begin{proof}
By Proposition \ref{prop:dominant-label}, there is a dominant class $C\subseteq Q$ with label $m$. For $p,q\in C$, the centred lift is exactly $m$. Hence
\[
\mathcal Q(a)
\ge
m^2\sum_{p<q\in C}\frac1{p^2q^2}
\ge
c_\rho m^2\sigma_Q^2,
\]
which proves \eqref{eq:4.24}.

Let $E_0=Q\setminus C$. The label bound and the low-energy hypothesis imply, for large $Y$,
\[
|m|\le\frac{|C|Y}{128}.
\tag{4.25}\label{eq:4.25}
\]
For each $q\in E_0$, the residue difference $d_q=a_q-m$ is nonzero modulo $q$. At most one prime in $C$ can divide an integer representative of $d_q$; applying Lemma \ref{lem:reciprocal-dispersion} to the remaining primes of $C$ and arguing as in \eqref{eq:4.8} gives
\[
\sum_{p\in C}
\left(\frac{H_{pq}}{pq}\right)^2
\ge
c_\rho\frac{|C|^3}{Y^2}
\tag{4.26}\label{eq:4.26}
\]
for each $q\in E_0$. Summing over $q\in E_0$ yields
\[
\mathcal Q(a)
\ge
c_\rho |E_0|\frac{N^3}{Y^2}.
\tag{4.27}\label{eq:4.27}
\]
Since $N\asymp Y/\log Y$, the low-energy hypothesis \eqref{eq:4.9} forces $|E_0|<1$ once $c_w$ is sufficiently small. Thus $E_0=\varnothing$, proving \eqref{eq:4.23}.
\end{proof}

\subsection{The fingerprint}

For $a$ with $\mathcal Q(a)\ge c_wY/\log^3Y$, put
\[
R=\mathcal Q(a),
\qquad
\eta(R)=
\min\left(\frac1{16},\frac{c_0R\log^3Y}{Y}\right),
\tag{4.28}\label{eq:4.28}
\]
where $c_0>0$ is a sufficiently small absolute constant. Choose an arbitrary base prime $p_0\in Q$. Call an edge $\{p_0,q\}$ fingerprint-good if
\[
|H_{p_0q}|
\le
B(R):=A\sqrt{\frac{R}{\eta(R)N}}\,Y^2,
\tag{4.29}\label{eq:4.29}
\]
with $A$ a large absolute constant. By Markov's inequality, all but at most $\eta(R)N$ neighbours are fingerprint-good once $A$ is large enough.

For each good neighbour $q$, the integer $H_{p_0q}$ determines $a_q\pmod q$. The fingerprint records $a_{p_0}$ and all these integers. Hence it determines the residues on at least $(1-\eta(R))N$ coordinates.

For two distinct fingerprint-good neighbours $q,q'$ with corresponding integers $w,w'$, suppose both residues $w,w'$ represented the same candidate at $p_0$ while remaining within the energy allowance. If $w\ne w'$, the reciprocal-dispersion bound \eqref{eq:4.5} applied to a retained set $F$ of $s\ge16$ good neighbours gives
\[
c_D\frac{s^3}{Y^2}
\le
\sum_{p\in F}
\left\|\frac{(w-w')p^{-1}}{q}\right\|^2
\le
2t_q(w)+2t_q(w')
<4g\frac{s^3}{Y^2},
\]
contradicting the choice $4g<c_D$. Thus the recorded values are rigid: two distinct candidate residues cannot both lie below the fingerprint energy threshold. This is the quantitative fingerprint contradiction used below.

The number of possible integer records on one good edge is at most
\[
1+\frac{4B(R)}{Y}
\le
C\left(1+\sqrt{\frac{RY}{\eta(R)N}}\right).
\tag{4.30}\label{eq:4.30}
\]
The omitted coordinates contribute at most
\[
\prod_{p\in E}p
\le(2Y)^{\eta(R)N}
\tag{4.31}\label{eq:4.31}
\]
possible completions. Therefore the logarithm of the number of completions of one fingerprint is
\[
O\bigl(\eta(R)N\log Y\bigr),
\tag{4.32}\label{eq:4.32}
\]
while the Bernoulli modulus contributes $e^{-\kappa R}$.

At the lower forcing scale $R\asymp Y/\log^3Y$, one has $\eta(R)\asymp R\log^3Y/Y$, so
\[
\eta(R)N\log Y
=O\!\left(\frac{R}{\log Y}\right)
=o(R).
\tag{4.33}\label{eq:4.33}
\]
For larger $R$, either the same estimate remains available or $R$ is already large enough that the trivial number of assignments is dominated by $e^{-\kappa R}$. In either case, after summing fingerprints, the total weighted mass at energy at least $R$ is $O(e^{-cR})$ for some $c>0$.

\begin{proposition}[Weighted noncoherent tail]\label{prop:noncoherent-tail}
There is $c>0$ such that
\[
\sum_{a:\,\mathcal Q(a)\ge c_wY/\log^3Y}
\prod_{p<q\in Q}
\left|(1-\theta)+\theta e\!\left(\frac{H_{pq}(a)}{pq}\right)\right|
\le
\exp\!\left(-c\frac{Y}{\log^3Y}\right).
\tag{4.34}\label{eq:4.34}
\]
\end{proposition}

\begin{proof}
The preceding fingerprint decomposition is summed dyadically in $R$. On each dyadic block the completion entropy is $o(R)$ at the forcing scale and remains absorbable thereafter, while Lemma \ref{lem:bernoulli-modulus} supplies $e^{-\kappa R}$. Summing the dyadic blocks gives \eqref{eq:4.34}.
\end{proof}

\subsection{Anchor partition}\label{subsec:anchor-partition}

Apply the preceding statements with
\[
Q=\mathcal B,
\qquad
Y=Z/2.
\]
Let
\[
F_{\mathcal B}=c_w\frac{Z}{\log^3Z}
\tag{4.35}\label{eq:4.35}
\]
be the anchor forcing floor. The coherent assignments are exactly those represented by one integer $m$ on every top coordinate. The noncoherent assignments have total anchor modulus
\[
O(e^{-cF_{\mathcal B}}).
\tag{4.36}\label{eq:4.36}
\]
For a coherent assignment, define
\[
\sigma_{\mathcal B,0}^2
=\sum_{p<q\in\mathcal B}\frac1{p^2q^2}.
\tag{4.37}\label{eq:4.37}
\]
Then
\[
\sigma_{\mathcal B,0}^2
\asymp\frac1{Z^2\log^2Z}.
\tag{4.38}\label{eq:4.38}
\]
The internal anchor modulus of label $m$ is at most
\[
\exp(-c m^2\sigma_{\mathcal B,0}^2).
\tag{4.39}\label{eq:4.39}
\]
The label bound \eqref{eq:4.24} and the forcing floor imply
\[
|m|
\ll
\frac{\sqrt{F_{\mathcal B}}}{\sigma_{\mathcal B,0}}
\ll
\frac{Z^{3/2}}{\sqrt{\log Z}}.
\tag{4.40}\label{eq:4.40}
\]
Finally, the weighted number of coherent labels is
\[
P_{\mathrm{top}}
:=
\sum_{m\ \mathrm{coherent}}
\exp(-c m^2\sigma_{\mathcal B,0}^2)
\ll
\sigma_{\mathcal B,0}^{-1}
\ll
Z\log Z.
\tag{4.41}\label{eq:4.41}
\]
This weighted partition, rather than a raw count of top CRT assignments, is the input to the row-compression argument.